\documentclass[10pt,a4paper]{amsart}
\usepackage[margin=19mm]{geometry}
\usepackage{amsmath,amssymb,mathtools}
\usepackage[scr=rsfs,cal=euler]{mathalfa}
\usepackage{xfrac}
\usepackage[unicode,hidelinks,bookmarks=false]{hyperref}
\newcommand{\ie}{i.e.\ }
\newcommand{\cf}{cf.\ }
\newcommand{\resp}{respectively\ }
\newcommand{\Subsection}{\subsection}
\providecommand{\nobreakdash}{\nobreak\hbox{-}}
\setlength{\emergencystretch}{2em}
\multlinegap0pt
\usepackage{enumitem}
\newcommand{\RS}{\operatorname{RS}}
\newcommand{\cB}{\mathcal B}
\newcommand{\supp}{\operatorname{supp}}
\theoremstyle{plain}
\newtheorem{theorem}{Theorem}[section]
\newtheorem{lemma}[theorem]{Lemma}
\newtheorem{proposition}[theorem]{Proposition}
\newtheorem{corollary}[theorem]{Corollary}
\theoremstyle{definition}
\newtheorem{definition}[theorem]{Definition}
\newtheorem{example}[theorem]{Example}
\theoremstyle{remark}
\newtheorem{remark}[theorem]{Remark}
\title[Cubical Sheaf Complexes with Constant Expansion with Applications to Asymptotically Good qLTCs]{Cubical Sheaf Complexes with Constant Expansion with Applications to Asymptotically Good qLTCs}
\author{Yeyuan Chen, Miryam Mi-Ying Huang, Yinchen Liu, Er-Cheng Tang}
\date{Source: arXiv:2609.28028v1 (CC BY 4.0)}
\begin{document}
\maketitle

\noindent\emph{Source and licence.} The delivered work is the article shown in the title above, version arXiv:2609.28028v1 (CC BY 4.0), distributed under the Creative Commons Attribution 4.0 International licence, as recorded in the arXiv licence metadata for this exact version and shown on the abstract page saved with this item. The full licence notice, which carries the licence URL and the address of that abstract page, is the bundled file \texttt{licenses/arxiv-2609.28028-CC-BY-4.0.txt}.
\par\smallskip

\noindent\textit{Editorial scope.} Counted unit: the article's Theorem (source label thm:inner-generation-main) (source0.tex lines 2203--2216, source label \texttt{thm:inner-generation-main}): ``Inner generation of polynomial zero sets For every , there exists a constant such that the following holds for every non-empty set . itemsep 0pt Let be chosen as in , be any non-empty set, and for eac''. It is delivered together with the article's complete original proof (source0.tex lines 3149--3245, one proof environment printed later in the article, detached from the statement by the article's own arrangement, 97 lines), copied line for line from the article's own text. Required context retained uncounted: eq:field-size-choice (lines 1908--1911); eq:tensor-parity (lines 2871--2876); lem:finite-union (lines 2893--2899); arb:axis-crt (lines 2915--2930); arb:rank (lines 2957--2961); arb:cylinders (lines 2989--3000). Editorial formatting only: an AMS \texttt{amsart} preamble that restates the article's own macro definitions and its own theorem declarations, and the theorem environments the retained lines use; the packages enumitem are loaded to match the article's own setup; the article's own class file is neither shipped nor input; the retained environments are renumbered sequentially in order of appearance, whereas the article prints them with its own numbering; the article's equation numbering is not reproduced and every cross-reference inside the retained text resolves to the wrapper's numbering. No citation occurs inside any retained line, so the wrapper carries no bibliography; All retained bodies are exact copies of the bundled \texttt{sources/211546/source0.tex}, so no omission receipt applies. No asset-inclusion command occurs inside any retained span and the package ships no image asset.


\section*{Retained context: eq:field-size-choice (source0.tex lines 1908--1911)}

\begin{equation}
 a_i=2(A+i),\qquad Q_i=2^{A+i},\qquad q_i=Q_i^2.
 \label{eq:field-size-choice}
\end{equation}


\section*{Retained context: eq:tensor-parity (source0.tex lines 2871--2876)}

\begin{equation}
 M\in\cB_h\quad\Longleftrightarrow\quad
 \sum_x M(x)f(x)\prod_i v_i(x_i)=0
 \quad\text{for every }f\text{ with }\deg_i f<h_i.
 \label{eq:tensor-parity}
\end{equation}


\section*{Retained context: lem:finite-union (source0.tex lines 2893--2899)}

\begin{lemma}[Unions and separation]\label{lem:finite-union}
If $W_\nu$ admits interpolation with bound $w_\nu$, then
$\bigcup_\nu W_\nu$ admits interpolation with bound $\sum_\nu w_\nu$.
If $z\notin W$ and $W$ admits interpolation with bound $w$, some
polynomial of coordinate degrees at most $w$ vanishes on $W$
but not at $z$.
\end{lemma}


\section*{Retained context: arb:axis-crt (source0.tex lines 2915--2930)}

\begin{lemma}[Isolating a derivative]\label{arb:axis-crt}\label{lem:CRT}
There exists a constant $K''_r \ge 1$ depending only on $r$ with the following property.
Suppose no variable divides $F$, and write
$F=P((X_i^{2^{u_i}})_i)$, extracting the largest possible power
of two from the exponents of each variable. Put $u_i=0$ for
variables absent from $F$.
For each variable $X_j$ on which $F$ depends, there is a polynomial
$R_j$ with the same zero set as $F$ on the grid $\prod_i S_i$ with $\deg_iR_j\le K''_r\deg_iF$.
At every grid point $x \in \prod_i S_i$ with $F(x)=0$, we have $\partial_iR_j(x)=0$ for $i\ne j$, and
$\partial_j R_j(x)$ is non-zero if and only if
\[
 b_j=(\partial_jP)((X_i^{2^{u_i}})_i)
\]
is nonzero on $x$. The polynomial $b_j$ is nonzero, has smaller total
degree than $F$, and has no larger coordinate degrees.
\end{lemma}


\section*{Retained context: arb:rank (source0.tex lines 2957--2961)}

\begin{lemma}\label{arb:rank}\label{lem:rank}
Let $M\in\cB_h$ on an $m$-dimensional grid, and let $S=\supp M$.
If $H\ge mB$, no $m$ polynomials of coordinate degrees at most $B$
vanishing on $S$ have independent gradients at a point of $S$.
\end{lemma}


\section*{Retained context: arb:cylinders (source0.tex lines 2989--3000)}

\begin{theorem}[Cylinder unions]
\label{arb:cylinders}\label{thm:cylinders}
For each nonempty $I\subseteq[m]$, let
$W_I\subseteq\prod_{i\in I}S_i$ admit interpolation with bound $w_I$.
Omit empty bases and put
\[
 Z=\bigcup_I\left(W_I\times\prod_{j\notin I}S_j\right),
 \qquad R=\max\{2,\sum_Iw_I\}.
\]
If $H\ge2mR$, then $\cB_h\cap\overline F^Z$ is generated by the
free-direction line codes in these cylinders.
\end{theorem}


\section{The counted unit: Theorem (source label thm:inner-generation-main) and its complete original proof}

\begin{theorem}[Inner generation of polynomial zero sets]
\label{thm:inner-generation-main}
For every $r \ge 2$, there exists a constant $K_r \ge 2$
such that the following holds for every non-empty set $J \subseteq [r]$.
\begin{itemize}[itemsep=0pt]
    \item Let $F, \Sigma_i$ be chosen as in \eqref{eq:field-size-choice}, $S_i \subseteq \Sigma_i$ be any non-empty set, and $n_i = |S_i|$ for each $i \in J$.
    \item Let $C_i = \RS(S_i,n_i-h_i)$ be a code with redundancy $h_i$ defined over $\overline{F}$ for $i \in J$.
    \item Let $E \in \overline{F}[X_i : i\in J]$ be a nonzero polynomial of degree at most $d \in \mathbb N$ in each variable.
\end{itemize}
If $h_i \ge K_r d$ for every index $i\in J$, then every codeword of
$\mathcal{B}((C_j)_{j\in J})$ supported on the zero set $Z_E$
is a sum of line-codewords whose associated
coordinate lines are contained in $Z_E$.
\end{theorem}

\begin{proof}[Proof of Theorem~\ref{thm:inner-generation-main}]
Fix $r \ge 2$. We induct on the number $m = |J|$ of chosen directions, simultaneously
over all direction subsets and all $S_i\subseteq\Sigma_i$.
Write $J=\{j_1,\ldots,j_m\}$ and use the same constant $K''_r$
from Lemma~\ref{lem:CRT} for every subset $J$.
For $m=1$, a nonzero polynomial of degree $d$ has at most $d$
roots, and the parity Vandermonde matrix is injective on them
when $H\ge2d$. Take $K'_1=2$.
Suppose the result is known in fewer than $m$ directions with
a common constant $K'_{m-1}$. We choose $K'_m = m^2 2^{m+1}(K'_{m-1}+K''_r+1)$, and prove the result for $m$ directions as follows.

\paragraph{Interpolation on bases.}
Fix a nonzero polynomial \(F\) of coordinate degrees at most \(d\).
Group the flats contained in \(Z_F\) that are maximal under
freeing coordinates by their fixed-coordinate sets \(I\), and let
\(W_I\) denote their sets of base points. Thus
$Z_F=\bigcup_I(W_I\times\prod_{j \in J \setminus I}S_j)$.
Fix a nonempty set \(I\subseteq J\), and write $F$ as a polynomial \(F(X_I,X_{J \setminus I})=\sum_\beta f_\beta(X_I)X_{J \setminus I}^{\beta}\) with free variables $X_{J \setminus I}$. 
Since $|S_i|\ge h_i\ge H\ge K'_m d>d$, a flat lies in $Z_F$
exactly when all these coefficients vanish at its base.
Choose a linear combination \(F_I=\sum_\beta\lambda_\beta f_\beta\) whose zero set $Z_{F_I}$ concides with the common zero set of $\{f_\beta\}_\beta$ on the grid \(S_I = \prod_{i \in I} S_i\). Such a choice exists over the infinite field $\overline{F}$ by avoiding the finitely many evaluation hyperplanes associated with points outside the common zero set.
By maximality, we have  $W_I=Z_{F_I}\setminus U_{F_I}$.
When $|I|<m$, since $F_I$ has coordinate degrees at most $d$, the induction hypothesis gives inner generation of $Z_{F_I}$ at redundancy $K'_{m-1}d$, which then implies that $W_I=Z_{F_I}\setminus U_{F_I}$ admits interpolation with bound $K'_{m-1}d$ by \eqref{eq:tensor-parity}.

We also need the following consequence: If such an $F$ has fewer than
$m$ variables, then any set
$T\subseteq Z_E\setminus U_E$ annihilated by $F$ admits interpolation
with bound $2^mK'_{m-1}d$.
It suffices to consider non-constant $F$. Use the
preceding base families in the variables appearing in $F$,
leaving all other directions free. Fix one such family
$W_I\times \prod_{j\in J \setminus I} S_j$, where $1\le|I|<m$. The preceding argument gives delta
interpolants $L_a(X_I)$ on $W_I$, with coordinate degrees
less than $K'_{m-1}d$.
For each $a\in W_I$, let
$T_a=\{y\in \prod_{j \in J\setminus I} S_j :(a,y)\in T\}$.
If $E(a,X_{J\setminus I})\equiv0$, the whole flat $\{a\}\times \prod_{j \in J\setminus I}S_{j}$
lies in $U_E$, so $T_a=\varnothing$. Otherwise, $E(a,X_{J\setminus I})$ is nonzero and has
coordinate degrees at most $d$, and $T_a\subseteq Z_{E(a,\cdot)}\setminus U_{E(a,\cdot)}.$
Induction and \eqref{eq:tensor-parity}, together with
interpolation passing to subsets, give delta interpolants
$L'_{a,b}(X_{J\setminus I})$ on each nonempty $T_a$, again with coordinate
degrees less than $K'_{m-1}d$.
The product $L_a(X_I)L'_{a,b}(X_{J\setminus I})$ is then a delta interpolant
for $(a,b)$ on $T\cap(W_I\times \prod_{j \in J\setminus I} S_j)$. Since the two factors
use disjoint variable sets, the coordinate-degree bound
remains unchanged. These sets cover $T$, so
Lemma~\ref{lem:finite-union} combines the fewer than $2^m$
families to give interpolation with bound $2^mK'_{m-1}d$.

\paragraph{Restricting the support to $U_E$.}
Let $M\in\mathcal{B}_h$ be supported on $Z_E$, with $H=\min_i h_i\ge K'_m d$. We also have $M\in\mathcal B_H$, where a scalar subscript means the same redundancy in every direction. We want to show that $\supp M\subseteq U_E$.
For each direction $i$, write $E$ as a polynomial in $X_i$ and take a
linear combination of its coefficients that vanishes exactly on the bases of
the direction-$i$ lines contained in $Z_E$. Let $P_0$ be the product of these $m$ polynomials. Then $P_0$ has grid zero set exactly $U_E$ and coordinate degrees at most $(m-1)d$. Put $N=P_0M$. It suffices to prove $N=0$.

Suppose otherwise. Note that  $N = P_0 M \in \mathcal B_{H-(m-1)d}$ with \(H-(m-1)d\ge m(K''_r+m)d\) by our choice of $K'_m$. We will reach a contradiction using  Lemma~\ref{lem:rank}. Starting with $Q_0=1$, we will recursively construct polynomials $Q_t$ with each coordinate degree at most $td$ and $\supp(Q_t N)\ne\varnothing$.
Choose $F_t$ as a polynomial with minimum total degree among \[\{G \neq 0: G|_{\supp(Q_t N)} = 0\; \text{  and } \deg_i G \le d \text{ for all } i\} \supseteq \{E\}.\]
No variable can divide $F_t$ because of its minimality and that all grid coordinates are nonzero. By the preceding
interpolation argument, $F_t$ uses all $m$ variables, provided
$H-2(m-1)d\ge2^mK'_{m-1}d$.
Applying Lemma~\ref{lem:CRT} to $F_t$ in direction $j_{t+1}$, we obtain $R_{t+1}, b_{t+1}$ with the following properties:
\begin{itemize}
    \item $R_{t+1} \neq 0$ has coordinate degrees at most $K''_r d$ and $R_{t+1}|_{\supp (Q_t N)} = 0$.
    \item $b_{t+1} \neq 0$ has coordinate degrees at most $d$ and total degree smaller than that of $F_t$.
    \item For every $x \in \supp(Q_t N)$, $\partial_j R_{t+1}(x) \neq 0$ if and only if $j = j_{t+1}$ and $b_{t+1}(x) \neq 0$.
\end{itemize}
Set $Q_{t+1}=Q_tb_{t+1}$. Minimality of $F_t$ implies that
$b_{t+1}$ does not vanish entirely on $\supp(Q_t N)$, and therefore $\supp(Q_{t+1} N) \neq \varnothing$.
Construct $A_{t+1} = Q_t R_{t+1}$, which satisfies the following:
\begin{itemize}
    \item $A_{t+1}$ has coordinate degrees at most
$(K''_r+t)d$.
    \item $A_{t+1}|_{\supp(N)} = 0$, since $R_{t+1}|_{\supp(Q_t N)} = 0$ and $Q_t|_{\supp(N)\setminus \supp(Q_t N)} = 0$.
    \item On $\supp(Q_{t+1} N) \subseteq \supp(Q_{t} N)$, both $Q_t$ and $b_{t+1}$ are nowhere zero, and $R_{t+1} = 0$. Hence $\nabla A_{t+1}=Q_t\nabla R_{t+1}$ has exactly one
nonzero component, in direction $j_{t+1}$.
\end{itemize}
After $m$ steps the gradients of $A_1,\ldots,A_m$ are linearly independent
on $\supp(Q_{m} N)\ne\varnothing$, contradicting Lemma~\ref{lem:rank}.
Hence $N = 0$, and therefore $\supp M\subseteq U_E$.

\paragraph{Gluing the lines.}
Group the positive-dimensional cylinders in \(Z_E\) that are
maximal under freeing coordinates by their fixed-coordinate sets
\(\varnothing\ne I\subsetneq J\), and denote their base sets
by \(W_I\). Then
$U_E=\bigcup_{\varnothing\ne I\subsetneq J}
\left(W_I\times\prod_{j\in J \setminus I}S_j\right)$.
We know that \(M\) has support on \(U_E\). Each $W_I$ involved admit interpolation with bound $K'_{m-1}d$ by the
first paragraph. Since there are fewer than \(2^m\) indices $I$ and we have chosen
$H\ge K'_md\ge 2m\,2^mK'_{m-1}d$, Theorem~\ref{thm:cylinders} decomposes \(M\) into line-codewords
on complete coordinate lines contained in \(U_E \subseteq Z_E\). This completes the induction.

Finally, setting $K_r = \max(K'_1,\dots,K'_r)$ completes the proof.

%Finally, $\cB_h\cap\overline F^{Z_E}$ and the span of line codes contained in $Z_E$ are kernels or images of matrices over the coefficient field. Field extension preserves their ranks, so equality over an algebraic closure proves equality over that field and over every extension.
\end{proof}

\end{document}
